\chapter{Round Robin versus Random Allocation}\label{ch:random}

FUNCTION: This chapter compares round robin with allocating the goods somehow. It introduces
the balanced random allocation as a generous benchmark, shows where round robin is far ahead and
where it fails EFX systematically, and presents reversing the last round as a simple remedy that
is just as easy to carry out by hand.

\section{Balanced Random Allocation as a Generous Benchmark}\label{sec:random}
FUNCTION: Defines the balanced random allocation (same bundle sizes as round robin, values
ignored) and explains why it is a generous model of allocating without any method. Figure:
EFX rate of round robin and of random allocation against Spearman in the reference cell. With
little or moderate agreement, round robin is far ahead (measured).

\section{Where Round Robin Fails Systematically}\label{sec:random-failure}
FUNCTION: With full agreement and a number of goods not divisible by the number of agents, round
robin fails to reach EFX with certainty (proven, Section~\ref{sec:divisibility-identical}), while
random allocation fails only almost always. Wording: systematic failure to reach EFX, never
unfair, since round robin remains EF1. Figure: full agreement, $n=4$, EFX rate against $m$ for
round robin, last round reversed and random allocation.

\section{Reversing the Last Round as a Simple Remedy}\label{sec:reversed}
FUNCTION: Short and exploratory. Reversing the last round removes exactly the pairs at which
round robin fails there. The guarantees are results of other authors (Manurangsi and Suksompong
2021; Amanatidis, Markakis and Ntokos 2020). Measured: it helps substantially when agents receive
few goods or values are similar.
